\begin{figure}[htbp]
\centering
\textbf{Fourier Series Decomposition of a Square Wave}

\vspace{0.5cm}

\begin{tikzpicture}
\begin{axis}[
    width=13cm, height=5cm,
    xlabel={Time $t$},
    ylabel={Amplitude},
    xmin=-1.1, xmax=1.1,
    ymin=-1.5, ymax=1.5,
    xtick={-1, -0.5, 0, 0.5, 1},
    xticklabels={$-T$, $-T/2$, $0$, $T/2$, $T$},
    ytick={-1, 0, 1},
    grid=major,
    grid style={gray!30},
    legend pos=outer north east,
    legend style={font=\footnotesize},
    title={\normalsize Individual Harmonic Components},
    title style={at={(0.5,1.05)}, anchor=south}
]

% DC component (a_0 = 0 for symmetric square wave)

% Fundamental (n=1): (4/pi)*sin(pi*x)
\addplot[UofSCGarnet, very thick, domain=-1.1:1.1, samples=200]
    {(4/pi)*sin(deg(pi*x))};
\addlegendentry{$n=1$: $(4/\pi)\sin(\omega_0 t)$}

% 3rd harmonic: (4/3pi)*sin(3*pi*x)
\addplot[UofSCAtlantic, thick, domain=-1.1:1.1, samples=200]
    {(4/(3*pi))*sin(deg(3*pi*x))};
\addlegendentry{$n=3$: $(4/3\pi)\sin(3\omega_0 t)$}

% 5th harmonic: (4/5pi)*sin(5*pi*x)
\addplot[UofSCHorseshoe, thick, domain=-1.1:1.1, samples=200]
    {(4/(5*pi))*sin(deg(5*pi*x))};
\addlegendentry{$n=5$: $(4/5\pi)\sin(5\omega_0 t)$}

% 7th harmonic
\addplot[UofSCCongaree, thick, domain=-1.1:1.1, samples=200]
    {(4/(7*pi))*sin(deg(7*pi*x))};
\addlegendentry{$n=7$: $(4/7\pi)\sin(7\omega_0 t)$}

\end{axis}
\end{tikzpicture}

\vspace{0.8cm}

\begin{tikzpicture}
\begin{axis}[
    width=13cm, height=5cm,
    xlabel={Time $t$},
    ylabel={Amplitude},
    xmin=-1.1, xmax=1.1,
    ymin=-1.5, ymax=1.5,
    xtick={-1, -0.5, 0, 0.5, 1},
    xticklabels={$-T$, $-T/2$, $0$, $T/2$, $T$},
    ytick={-1, 0, 1},
    grid=major,
    grid style={gray!30},
    legend pos=outer north east,
    legend style={font=\footnotesize},
    title={\normalsize Cumulative Sum of Harmonics},
    title style={at={(0.5,1.05)}, anchor=south}
]

% Ideal square wave (for reference)
\addplot[UofSC50Black, dashed, very thick, domain=-1.1:-0.5, samples=2, forget plot] {-1};
\addplot[UofSC50Black, dashed, very thick, domain=-0.5:-0.5, samples=2, forget plot] coordinates {(-0.5,-1) (-0.5,1)};
\addplot[UofSC50Black, dashed, very thick, domain=-0.5:0.5, samples=2, forget plot] {1};
\addplot[UofSC50Black, dashed, very thick, domain=0.5:0.5, samples=2, forget plot] coordinates {(0.5,1) (0.5,-1)};
\addplot[UofSC50Black, dashed, very thick, domain=0.5:1.1, samples=2] {-1};
\addlegendentry{Ideal Square Wave}

% n=1 only
\addplot[UofSCGarnet, thick, domain=-1.1:1.1, samples=200]
    {(4/pi)*sin(deg(pi*x))};
\addlegendentry{$N=1$}

% n=1,3
\addplot[UofSCAtlantic, thick, domain=-1.1:1.1, samples=200]
    {(4/pi)*sin(deg(pi*x)) + (4/(3*pi))*sin(deg(3*pi*x))};
\addlegendentry{$N=3$}

% n=1,3,5
\addplot[UofSCHorseshoe, thick, domain=-1.1:1.1, samples=200]
    {(4/pi)*sin(deg(pi*x)) + (4/(3*pi))*sin(deg(3*pi*x)) + (4/(5*pi))*sin(deg(5*pi*x))};
\addlegendentry{$N=5$}

% n=1,3,5,7,9,11
\addplot[UofSCCongaree, very thick, domain=-1.1:1.1, samples=300]
    {(4/pi)*sin(deg(pi*x)) + (4/(3*pi))*sin(deg(3*pi*x)) + (4/(5*pi))*sin(deg(5*pi*x)) + (4/(7*pi))*sin(deg(7*pi*x)) + (4/(9*pi))*sin(deg(9*pi*x)) + (4/(11*pi))*sin(deg(11*pi*x))};
\addlegendentry{$N=11$}

% n=1,3,5,7,9,11,13,15,17,19
\addplot[UofSCRose, very thick, domain=-1.1:1.1, samples=400]
    {(4/pi)*sin(deg(pi*x)) + (4/(3*pi))*sin(deg(3*pi*x)) + (4/(5*pi))*sin(deg(5*pi*x)) + (4/(7*pi))*sin(deg(7*pi*x)) + (4/(9*pi))*sin(deg(9*pi*x)) + (4/(11*pi))*sin(deg(11*pi*x)) + (4/(13*pi))*sin(deg(13*pi*x)) + (4/(15*pi))*sin(deg(15*pi*x)) + (4/(17*pi))*sin(deg(17*pi*x)) + (4/(19*pi))*sin(deg(19*pi*x))};
\addlegendentry{$N=19$}

\end{axis}
\end{tikzpicture}

\caption{Fourier series decomposition of a square wave. \textbf{Top:} Individual odd harmonic components ($n = 1, 3, 5, 7$) with amplitudes $4/(n\pi)$. A symmetric square wave contains only odd harmonics. \textbf{Bottom:} Cumulative sum showing how adding more harmonics improves the approximation. The dashed gray line shows the ideal square wave. Note the Gibbs phenomenon (overshoot near discontinuities) which persists regardless of how many terms are included. The Fourier series formula is: $f(t) = \sum_{n=1,3,5,\ldots}^{\infty} \frac{4}{n\pi}\sin(n\omega_0 t)$ where $\omega_0 = 2\pi/T$.}
\label{fig:fourier_series}
\end{figure}
